\subsection{Segunda iteración de la Fase 4: Cambios necesarios}
\label{subsec:F4I2}

\subsubsection{Velocidades por plataforma}

\begin{itemize}
	\item El GWI gana el campo \texttt{pausas\{simulador, externo\}}: un
	solo archivo de cartucho define las dos velocidades de reloj. El
	simulador lee siempre, como es lógico, \texttt{simulador}; el firmware
	del ESP32 leerá \texttt{externo}. La lógica del motor de estados no
	cambia en absoluto: es determinista en \textit{ticks}, y la pausa solo
	decide a qué velocidad real se consumen.
	
	\item Si el campo no está presente en el \emph{setup.json} o en el
	cartucho, el Director usa por defecto 500 ms en ambas plataformas.
	
	\item La ficha del cartucho del diseñador gráfico (la que se abre al
	pulsar en una zona vacía) gana los dos campos, con validación de entero
	positivo. Los valores viajan en el historial (deshacer y rehacer) y en
	el GWI al guardar.
	
	\item El puente de datos (\texttt{comunListas}) gana los accesos
	\texttt{fijarPausas} y \texttt{pausasDeCartucho}, siguiendo el patrón de
	la Secuencia y del fin de partida: \texttt{baseDirector} los escribe al
	cargar el cartucho y el diseñador los sincroniza en vivo al final de cada
	reconstrucción de la ficha, de modo que cualquier edición es visible para
	el motor sin guardar ni recargar.
\end{itemize}

\subsubsection{Aceleración del juego}

\begin{itemize}
	\item El GWI gana el campo
	\texttt{aceleracion \{ cada, resta, minimo \}}: cada \texttt{cada}
	\textit{ticks} de partida, la pausa vigente se reduce en \texttt{resta}
	milisegundos, sin bajar nunca de \texttt{minimo}.
	
	\item Un solo conjunto de parámetros sirve a ambas plataformas: la
	aceleración se aplica sobre la pausa inicial de cada una, por lo que el
	efecto es proporcional y no hace falta duplicar los campos.
	
	\item Si el campo no está presente, la pausa permanece fija durante toda
	la partida, por tanto la velocidad será constante.
	
	\item El Director aplica la reducción en su contador de \textit{ticks} y
	expone la pausa vigente, que la barra de estado con \texttt{GW\_DEBUG}
	presenta junto al número de \textit{tick}, para poder verificar la
	aceleración durante las pruebas.
	
	\item La aceleración cuenta \textit{ticks} de partida, no de reloj: el
	motor de estados sigue siendo determinista, de modo que dos partidas con
	la misma semilla atraviesan exactamente los mismos estados, solo que a
	velocidad creciente.
\end{itemize}

\subsubsection{Política de lanzamiento de rutas}

\begin{itemize}
	\item El GWI gana el campo \texttt{Nº para aumentar rutas}
	(clave \texttt{rutas.aumentarCada}).
	
	\item El Director inicia cada partida con el \textbf{total de rutas
		lanzadas} en 0 y el \textbf{máximo de rutas simultáneas} en 1; ambos
	valores se reinician con cada partida.
	
	\item Cada vez que se lanza una nueva ruta, el \textbf{total de rutas
		lanzadas} aumenta en 1; cuando una ruta muere, este total no cambia.
	
	\item Cuando el \textbf{total de rutas lanzadas} es múltiplo del campo
	\texttt{Nº para aumentar rutas}, el \textbf{máximo de rutas
		simultáneas} aumenta en 1, sin superar nunca el número de rutas
	definidas en el cartucho.
	
	\item El Director lleva además el \textbf{contador de rutas vivas}: el
	número de rutas en pantalla en cada momento. En cada \textit{tick}, si
	las rutas vivas son menos que el máximo de simultáneas, se lanza una
	nueva ruta. La regla actual (nacer cuando no queda ninguna viva) es el
	caso particular de \emph{máximo de rutas simultáneas = 1}.
\end{itemize}

\subsubsection{Combinación de pulsaciones del teclado}
Este era un tema pendiente, se creó para permitir combinar teclas para ejecutar saltos. La regla \texttt{combinar} decide cómo se procesan las pulsaciones encoladas del teclado: con ella activa, el buffer se reduce a ejes y el jugador da un paso por eje (dos pulsaciones iguales en la misma ventana cuentan como una, y las opuestas se anulan, permitiendo el movimiento diagonal); sin ella, cada pulsación es un paso independiente por orden de llegada. Hoy esta regla vive únicamente en el \emph{setup}, fuera del cartucho, de modo que no es editable desde el diseñador gráfico ni viaja en el GWI: un cartucho definitivo no puede reproducir el comportamiento de su teclado.

Pero es un tema que solo usaban algunos cartuchos de última generación, por lo que finalmente \underline{la decisión ha sido dejarlo como posible ampliación del proyecto}.

\subsubsection{Validación de los cambios}

\begin{itemize}
	\item \textbf{Contador de rutas vivas atascado en la decisión inmediata
		de colisión}: Al probar la nueva política de rutas con un máximo de
	simultáneas mayor que 1, el contador de rutas vivas se quedaba
	atascado en 1: solo se producía una muerte de ruta y, a partir de ahí,
	el motor dejaba de lanzar rutas nuevas aunque el máximo lo permitía.
	La causa era que el decremento del contador al morir una ruta solo
	se aplicaba en cuatro de los cinco puntos del motor donde una ruta
	muere; faltaba el de la decisión inmediata al avanzar a un punto de
	colisión (la ruta muere \emph{en el momento} de chocar, sin esperar al
	tick siguiente). Con el decremento llamado también desde ese punto, el
	contador baja en todas las muertes y el lanzamiento vuelve a
	funcionar. Este error es invisible con la regla histórica (máximo de
	simultáneas = 1), porque la ruta nueva nace justamente cuando no queda
	ninguna viva: por eso no apareció hasta que la fase dio al cartucho el
	control de la política de lanzamiento.
	
	\item \textbf{Error de escritura en el nombre del puente.} Durante el
	desarrollo, una línea del \texttt{juegoDirector} llamaba al puente con
	una doble eme (\texttt{CommunListas}), lo que en el navegador habría
	provocado un fallo silencioso de la aceleración (el Director no
	encuentra el campo y aplica la pausa fija). Cazado en la simulación
	antes de la prueba en navegador.
\end{itemize}